% Vietnamese translation for OI Public Library
% Source-File: IOI中国国家候选队论文/2007/day2/2.古楠《平面嵌入》.ppt
\documentclass[11pt]{article}
\usepackage{oipl-vn}

\newcommand{\Oh}{\mathcal{O}}
\newcommand{\Kfive}{K_5}
\newcommand{\Kthree}{K_{3,3}}

\title{Bản trình chiếu: Nhúng phẳng}
\author{Gu Nan}
\date{2007}

\begin{document}
\maketitle

\origtitle{Slide companion for planar embedding}
\sourcepath{IOI China candidate team papers / 2007 / day 2 / Gu Nan.ppt}

\section{Phạm vi bản dịch}

Tệp PowerPoint đi kèm bài viết đã được dịch từ PDF về thuật toán nhúng phẳng.
Bản ghi chú này tóm lược các khái niệm chính, mục tiêu thuật toán và những
thành phần dữ liệu được slide nhắc tới.

\section{Khái niệm}

Một đồ thị là phẳng nếu có thể vẽ trên mặt phẳng sao cho các cạnh chỉ giao
nhau tại đỉnh chung. Một nhúng phẳng lưu thông tin này theo dạng tổ hợp: với
mỗi đỉnh, danh sách kề được sắp theo thứ tự quanh đỉnh.

Hai chướng ngại cổ điển là \(\Kfive\) và \(\Kthree\). Theo định lý Kuratowski,
một đồ thị không phẳng khi và chỉ khi chứa một đồ thị con đồng phôi với một
trong hai đồ thị ấy.

\section{Mục tiêu thuật toán}

Với đồ thị \(G\), thuật toán cần phán đoán tính phẳng và nếu phẳng thì dựng
một nhúng phẳng. Bài viết tập trung vào cách xử lý theo DFS, tách thành phần
song liên thông, cạnh cây, cạnh hồi và các thao tác nhúng cạnh.

\section{Nhúng cạnh}

Khi thêm cạnh vào đồ thị đang nhúng, ta phải bảo đảm cạnh mới nằm trên một
mặt phù hợp. Các slide nhắc tới hoạt động ngoài, liên quan và hoạt động trong:
đây là các nhãn giúp xác định thành phần nào có thể bị đảo, thành phần nào
cần ghép, và khi nào quá trình nhúng thất bại.

\section{Kết luận}

Nhúng phẳng là ví dụ tốt cho việc biến một khái niệm hình học thành cấu trúc
tổ hợp. Cài đặt khó nằm ở duy trì thứ tự quanh đỉnh và các thành phần song
liên thông, nhưng tư tưởng chính vẫn là thêm cạnh theo thứ tự kiểm soát được
và luôn giữ mọi ràng buộc của mặt ngoài.

\end{document}
